\subsection{Kirchhoffsche Gesetze 3}
\Aufgabe{
    Für die Brückenschaltung sind folgende Werte gegeben: $R_1 = 4\ \mathrm{k}\Omega, R_2=20 \ \mathrm{k}\Omega, R_3 = 6 \ \mathrm{k}\Omega, R_4 = 20 \ \mathrm{k}\Omega$. \newline
    Der Innenwiderstand des Voltmeters kann als unendlich groß angenommen werden. Berechnen Sie die folgenden Größen:
    \begin{enumerate}[label=\alph*)]
        \item  $I_{13}$ und $I_{24}$
        \item $U_1,\ U_2, \ U_3$ und $U_4$
        \item Die Spannung zwischen den Punkten A und B.
        \item Wann wäre die Brücke abgeglichen, d. h. wann ist die Spannung zwischen den Punkten A und B Null?
        Stellen Sie die Abgleichbedingung in möglichst allgemeiner Form als Funktion der Widerstände $R_1$ bis $R_4$ dar.      
    \end{enumerate}
    \begin{center}
        \input{Tikz/src/Aufgaben/Kirchhoff_A3}
    \end{center}
}

\Loesung{
	    Für die Brückenschaltung sind folgende Werte gegeben: $R_1 = 4\ \mathrm{k}\Omega, R_2=20 \ \mathrm{k}\Omega, R_3 = 6 \ \mathrm{k}\Omega, R_4 = 20 \ \mathrm{k}\Omega$. \newline
   \begin{center}
        \input{Tikz/src/Aufgaben/Kirchhoff_A3}
   \end{center}
   
   
        Der Innenwiderstand des Voltmeters kann als unendlich groß angenommen werden. Berechnen Sie die folgenden Größen:
    \begin{enumerate}[label=\alph*)]
        \item  $I_{13}$ und $I_{24}$
        
        Geringfügiges Umzeichnen erhöht die Übersicht:

        \begin{center}
            \begin{tikzpicture}
    \draw(0,6)to[short,i,name=zu,o-](4,6)
    to[short,-*](4,5.5)
    to[short](5.5,5.5)
    to [R,v,i_,l=$R_1$,name=R1](5.5,3.5)
    to[short,-*](5.5,3)
    to[short](5.5,2.5)
    to [R,v,i,l=$R_3$,name=R3](5.5,0.5)
    to[short,-*](4,0.5)
    to[short](4,0)
    to[short,-o](0,0);
    \draw(4,5.5)to[short](2.5,5.5)
    to [R,v,i_,l=$R_2$,name=R2](2.5,3.5)
    to[short,-*](2.5,3)
    to[short](2.5,2.5)
    to [R,v,i,l=$R_4$,name=R4](2.5,0.5)
    to[short,-*](4,0.5);


   % to [R,-*](2,1)
   % to [R,i<^,name=R2](0,3);
   % \draw(4,3) to[short,-o](4,0);
   % \draw(2,5)to[rmeterwa,t=V](2,1);

   % \iarrmore{zu}{$I$};
    \iarrmore{R1}{$I_{13}$};
    \varrmore{R1}{$U_1$};
   % \iarrmore{R2}{$I_{24}$};

    \node at (0,-0.4) {+};
    \node at (4,-0.4) {-};
    \node at (5.8,3) {A};
    \node at (1.7,3) {B};
   
    \draw[->] (0,5.7) -- (0,0.3) node[midway, right] {24 V};
\end{tikzpicture}
        \end{center}
        

        Zwei Möglichkeiten, um Teilströme zu ermitteln: Stromteiler oder Ohm´sches Gesetz.\\

        1. Möglichkeit: Stromteiler\\

        Gesamtstrom $I_0$ ermitteln:\\
        $I_0 = \frac{U_0}{R_\mathrm{ges}}$\\
        $R_\mathrm{ges}= (R_1 + R_3)||(R_2+R_4) = \frac{(R_1+R_3) \cdot (R_2+R_4)}{R_1 + R_2 + R_3 + R_4} = \frac{10 \, \mathrm{k} \Omega \cdot 40 \, \mathrm{k} \Omega}{50 \, \mathrm{k} \Omega} = 8 \, \mathrm{k} \Omega$\\
        $I_\mathrm{0} = \frac{24 \, \mathrm{V}}{8 \, \mathrm{k} \Omega} = 3 \, \mathrm{mA}$\\

        Teilstrom ausrechnen:\\
        $I_{13} = I_0 \cdot \frac{R_2+R_4}{R_1+R_2+R_3+R_4} = 3 \, \mathrm{mA} \cdot \frac{40 \, \mathrm{k} \Omega}{50 \, \mathrm{k} \Omega} = 2,4 \, \mathrm{mA}$\\

        Teilstrom $I_{24}$ mittels Knotenregel ausrechnen:\\
        $I_0 = I_{13} + I_{24} \rightarrow I_{24} = I_0 - I_{13} = 3 \, \mathrm{mA}- 2,4 \, \mathrm{mA} = 0,6 \, \mathrm{mA}$\\

        2. Möglichkeit, die hier, aber nicht immer anwendbar ist: Ohm´sches Gesetz:\\

        $I_{13} = \frac{U_0}{R1+R_3} = \frac{24 \, \mathrm{V}}{10 \, \mathrm{k} \Omega} = 2,4 \, \mathrm{mA}$\\

        $I_{24} = \frac{U_0}{R2+R_4} = \frac{24 \, \mathrm{V}}{40 \, \mathrm{k} \Omega} = 0,6 \, \mathrm{mA}$\\





        \item $U_1,\ U_2, \ U_3$ und $U_4$\\
        
        2 Möglichkeiten: Spannungsteiler oder Ohm´sches Gestz:\\
        
        Spannungsteiler:
        \begin{center}
        \begin{tikzpicture}
    \draw(4,5.5)to[short, o-](5.5,5.5)
    to [R,v,i,l=$R_1$,name=R1](5.5,3.5)
    to[short,-*](5.5,3)
    to[short](5.5,2.5)
    to [R,v,i,l=$R_3$,name=R3](5.5,0.5)
    to[short,-o](4,0.5);
 



   % to [R,-*](2,1)
   % to [R,i<^,name=R2](0,3);
   % \draw(4,3) to[short,-o](4,0);
   % \draw(2,5)to[rmeterwa,t=V](2,1);

   % \iarrmore{zu}{$I$};
    \iarrmore{R1}{$I_{13}$};

    \varrmore{R1}{$U_1$};

    \varrmore{R3}{$U_3$};

   % \iarrmore{R2}{$I_{24}$};


  
    \draw[->] (4,5.2) -- (4,0.8) node[midway, right] {24 V};
\end{tikzpicture}
        \end{center}
        $U_3 = 24 \, \mathrm{V} \cdot \frac{R_3}{R1+R_3} = 14.4 \, \mathrm{V}$\\
        Alternativ: \\
        $U_3 = R_3 \cdot I_{13} = 14,4 \, \mathrm{V}$\\

        $U_1$ kann beispielsweise mittels Maschenumlauf ermittelt werden, Spannungsteiler oder Ohm´sches Gesetz wären aber auch möglich:\\
        $U_1 = U_0 - U_3 = 24 \, \mathrm{V} - 14,4 \, \mathrm{V} = 9,6 \, \mathrm{V}$\\

        $U_2 = R_3 \cdot I_{24} = 12 \, \mathrm{V}$\\
        $U_4 = R_3 \cdot I_{24} = 12 \, \mathrm{V}$\\






        

        \item Die Spannung zwischen den Punkten A und B.\\
        
        Zur Bestimmung der Spannung $U_\mathrm{AB}$ werden Maschen in das Ersatzschaltbild eingezeichnet:
        \begin{center}
        \begin{tikzpicture}
    \draw(0,6)to[short,i,name=zu,o-](4,6)
    to[short,-*](4,5.5)
    to[short](5.5,5.5)
    to [R,v,i,l=$R_1$,name=R1](5.5,3.5)
    to[short,-*](5.5,3)
    to[short](5.5,2.5)
    to [R,v,i,l=$R_3$,name=R3](5.5,0.5)
    to[short,-*](4,0.5)
    to[short](4,0)
    to[short,-o](0,0);
    \draw(4,5.5)to[short](2.5,5.5)
    to [R,v,i,l=$R_2$,name=R2](2.5,3.5)
    to[short,-*](2.5,3)
    to[short](2.5,2.5)
    to [R,v,i,l=$R_4$,name=R4](2.5,0.5)
    to[short,-*](4,0.5);


   % to [R,-*](2,1)
   % to [R,i<^,name=R2](0,3);
   % \draw(4,3) to[short,-o](4,0);
   % \draw(2,5)to[rmeterwa,t=V](2,1);

   % \iarrmore{zu}{$I$};
    \iarrmore{R1}{$I_{13}$};
    \iarrmore{R2}{$I_{24}$};
    \varrmore{R1}{$U_1$};
    \varrmore{R2}{$U_2$};
    \varrmore{R3}{$U_3$};
    \varrmore{R4}{$U_4$};
   % \iarrmore{R2}{$I_{24}$};



    \node at (5.8,3) {A};
    \node at (1.7,3) {B};
    \draw[->] (5.2,3) -- (2.8,3) node[midway, below] {$U_\mathrm{AB}$};
    \draw[->] (0,5.7) -- (0,0.3) node[midway, right] {24 V};
    \draw[->,draw = blue, shift={(5.25,4.3)}] (180:.7cm) arc (15:345:.6cm) node at(4,4){};
    \draw[->,draw = blue, shift={(5.25,1.8)}] (180:.7cm) arc (15:345:.6cm) node at(5,2){};
    \draw(4,3.8) node[label={above:{\textcolor{blue}{$M1$}}}] {};
    \draw(4,1.3) node[label={above:{\textcolor{blue}{$M2$}}}] {};
\end{tikzpicture}
        \end{center}
        $M_1: U_2 - U_\mathrm{AB} - U_ 1 = 0$\\
        $\rightarrow  U_\mathrm{AB} = U_2 - U_1 = 12 \, \mathrm{V} - 9,6 \, \mathrm{V} = 2,4 \, \mathrm{V}$\\

        Kontrolle mit Masche $M2$:\\

        $M2: U_\mathrm{AB} + U_4 - U_2 = 0$\\
        $\rightarrow U_\mathrm{AB} = U_3 - U_4 = 14,4 \, \mathrm{V} - 12 \, \mathrm{V} = 2,4 \mathrm{V}$\\

        \item Wann wäre die Brücke abgeglichen, d. h. wann ist die Spannung zwischen den Punkten A und B Null?
        Stellen Sie die Abgleichbedingung in möglichst allgemeiner Form als Funktion der Widerstände $R_1$ bis $R_4$ dar. \\

        Es muss also gelten: $U_\mathrm{AB} = 0 \, \mathrm{V}$\\

        Auf $M1$ bezogen bedeutet dies: $U_1 = U_2 \rightarrow I_{13} \cdot R_1 = I_{24} \cdot R_2$\\
        Umgestellt nach $I_{24}$ ergibt sich: $I_{24} = I_{13} \cdot \frac{R_1}{R_2}$\\

        Auf $M2$ bezogen bedeutet dies: $U_3 = U_4 \rightarrow I_{13} \cdot R_3 = I_{24} \cdot R_4$\\

        Durch Einsetzen des Ausdruckes für $I_{24}$ aus $M1$ in die Gleichungen von $M2$ lässt sich $I_{24}$ eliminieren und das Gleichungssystem lösen.\\

        $\rightarrow I_{13} \cdot R_3 = I_{13} \cdot \frac{R_1 \cdot R_4}{R_3}$\\

        $\rightarrow \frac{R_1}{R_3} = \frac{R_2}{R_4}$
 
    


        
        
    \end{enumerate}
    \begin{center}
        
   
   
\end{center}
}